\section{Dataset-as-Input：数据矩阵、mask 与预测契约}

要让整个数据集进入模型，首先必须把 supervised-learning notation 改写为一个统一矩阵。NPT 将 rows 视为 datapoints、columns 视为具有共享语义的 attributes；label 也占据某个 attribute column。Mask matrix 决定哪些 entries 被观察，哪些 entries 是模型应预测的未知量。

\lecturefigure{slide-13-entire-dataset-input.jpg}{第一要素：把 entire dataset，即全部 datapoints，作为模型输入。}{Stanford Online 官方视频 00:24:00--00:24:13。}

设
\[
X\in\mathbb{R}^{n\times d},\qquad M\in\{0,1\}^{n\times d}.
\]

\begin{itemize}
  \item $n$：当前 forward 中的 datapoint 数；小数据可为全数据集，大数据通常为 mini-batch；
  \item $d$：attributes 数，既包含 input features，也可包含 target columns；
  \item $X_{ij}$：第 $i$ 个 datapoint 的第 $j$ 个 attribute；
  \item $M_{ij}=1$：该 entry 被 mask，需要预测；$M_{ij}=0$：该 entry 可见。
\end{itemize}

\lecturefigure{slide-14-datapoint-attention.jpg}{第二要素：在 datapoints 之间应用 self-attention。}{Stanford Online 官方视频 00:24:14--00:24:27。}

在 test-time batch 中，attention 可以同时连接 train--train、test--test 与 train--test rows。这里的“train/test”是由 target visibility 决定的角色：训练 rows 可以暴露部分 labels，测试 rows 的 labels 必须全部遮挡。模型看到同一个矩阵，却通过 mask 知道哪些 entries 可以作为 context。

\lecturefigure{slide-15-masking-objective-overview.jpg}{第三要素：用 masking-based objective 联合学习 target 与 feature reconstruction。}{Stanford Online 官方视频 00:24:28--00:24:41。}

定义 masked 与 observed entries：
\[
X^{\mathcal{M}}=\{X_{ij}\mid M_{ij}=1\},\qquad
X^{\mathcal{O}}=\{X_{ij}\mid M_{ij}=0\}.
\]
NPT 学习
\[
p\!\left(X^{\mathcal{M}}\mid X^{\mathcal{O}},M\right),
\]
并输出 $\widehat X\in\mathbb{R}^{n\times d}$。这种表示把分类、回归、imputation、self-supervision 与 semi-supervision 都转化为“改变 mask 的位置”。

\begin{importantbox}{Mask 是任务接口}
同一 architecture 可以通过 $M$ 表达不同问题：遮住测试 rows 的 label column 得到 supervised prediction；随机遮住 feature entries 得到 imputation/self-supervision；遮住多个 target columns 得到 multi-target prediction。模型接口统一，不代表不同任务的数据分布与 loss 已自动统一。
\end{importantbox}

\subsection{完整 notation：为什么必须把 mask 也输入模型}

前面的三张图给出了 dataset、datapoint attention 与 masking objective，本小节把它们收束成一个不会产生歧义的输入接口。读完整 notation 图时应先分清“数值是什么”和“数值是否可见”：同一个零既可能是合法观测，也可能是缺失占位符，仅看 $X$ 无法分辨。把 $M$ 与 $X$ 一同嵌入，相当于显式告诉模型每个值的 epistemic status：它是已知 evidence，还是待预测 query；这也决定 loss 应落在哪些 entries 上。

\lecturefigure{slide-16-full-dataset-mask-notation.jpg}{完整输入契约：dataset $X$、mask $M$、目标 $p(X^{\mathcal{M}}\mid X^{\mathcal{O}})$。}{Stanford Online 官方视频 00:25:06--00:27:15。}

\begin{knowledgebox}{读图：三种颜色与问号}
每一行是 datapoint，每一列是语义固定的 attribute。问号不是“值为零”，而是 $M_{ij}=1$ 的未知 entry。模型需要同时利用同一行其他 features、其他行 features，以及允许暴露的其他行 targets，来重建问号位置。
\end{knowledgebox}

\begin{warningbox}{Test-label leakage 的底线}
Stochastic target masking 可以在训练 rows 上保留部分 labels，帮助模型学习跨行 lookup；测试 rows 的真实 labels 永远不能作为输入。若构造 batch 或 mask 时把 test label 暴露，漂亮结果只说明 pipeline 泄漏答案。
\end{warningbox}

\teachervoice{讲者把三要素先整体说出，再逐个展开：数据集作为输入、datapoint attention、masking objective。这个顺序提醒我们，不应只盯 architecture block；真正定义 NPT 的还有 input semantics 与 training task。}

\subsection{本章小结}

NPT 用 $(X,M)$ 统一表示 evidence 与 query。Rows 是 datapoints，columns 是共享语义 attributes，mask 决定预测任务。这个契约使 train/test points 可以共同进入 attention，同时要求严格保护 test targets。

